\section{Conclusion}
\label{sec:conclusion}

TIGER (Tension, Interaction, Game-Architecture, Elegance, Range) is a five-dimensional
framework for characterizing board game design that measures distinctiveness rather than
quality. Formalized from PCA-identified components of the within-game normalized TIGRIS
corpus, TIGER provides anchor-descriptor scoring at 0, 5, and 10 for each dimension and
correctly discriminates $\geq 90\%$ of expert-rated game pairs in a blind panel study.

The framework's primary utility is in design conversations where quality rankings are
insufficient: recommendation (finding games with a similar design character to a known
favorite), design critique (identifying under-developed dimensions in a prototype), and
market analysis (identifying under-served design regions). The companion papers TG-B2,
TG-C1, and TG-C2 demonstrate these applications empirically.

TIGER scoring rubrics, profiles for 20 canonical games, and the expert panel validation
dataset are available at [URL].
